\documentclass[12pt]{article}
\usepackage[utf8]{inputenc}
\usepackage[T1]{fontenc}
\usepackage{hyperref}
\title{Uru-Pa-In Language Resources for K-12 Curriculum Development}
\author{Compiled by Claude Sonnet 5 (claude-sonnet-5) for Andr\'as Kornai}
\date{}
\begin{document}
\maketitle

\section{Introduction}
Uru-Pa-In (ISO 639-3 urp) is spoken by roughly 200 people in Rond\^onia, Brazil. Sources disagree on its classification. Some place it in Tupian, while others treat it as unclassified or possibly linked to Chapacuran. It should not be confused with Uru-Eu-Wau-Wau, which is a Kagwahiva variety. Searches found only general reference pages and no dedicated teaching or documentation materials.

\section{Primary Educational Resources}
No primers, textbooks or school programmes for Uru-Pa-In were found in search results.

\section{Linguistic and Reference Materials}
\begin{itemize}
\item Wikipedia entry on the language: \url{https://en.wikipedia.org/wiki/Uru-Pa-In_language}
\item A Portuguese-language article in the journal Lingua Viva (UNIR) that mentions the Uru Pa In, Amondawa and Uru Eu Wau Wau peoples. Its claim that they share one language conflicts with other sources and was not verified: \url{https://periodicos.unir.br/index.php/linguaviva/article/download/1133/1229/4102}
\end{itemize}
No dictionary or grammar was found.

\section{Religious and Community Resources}
Joshua Project has a people-group profile, which is not a community-produced resource: \url{https://joshuaproject.net/people_groups/15740/BR}

\section{Digital and Audio Resources}
No audio, video or corpus resources were found.

\section{Curriculum Development Context}
Almost nothing is publicly documented. Any curriculum would have to begin with community-led documentation, working with the Uru-Pa-In people and Brazilian indigenous education bodies. Contact details were not found and are not given here.
\end{document}